Replacing guided proposals and rollouts by uniform actions, while retaining
the leaf estimate, increases $L$ by $39.3\pm4.9$ steps across the four-task
seed blocks. On Four~3D and Narrow, success falls to $0/20$, compared with
$17/20$ and $16/20$ for the full variant. Domain-informed action selection
and rollout are therefore the strongest supported components in this study.

The other comparisons do not establish an improvement from the full
configuration. Removing leaf values changes $L$ by $-3.2\pm4.7$;
removing widening gives $-2.7\pm4.6$; guided control without search gives
$-2.2\pm4.9$. All three intervals include zero. The negative point estimates
favour simpler variants, but the data do not resolve their differences.
Widening controls storage and exploration allocation; the short budget already
limits growth. The leaf estimate ignores coupled occupancy and future slip
risk, so it can also bias finite-budget decisions. These results support
guidance, while leaving the benefit of these additional approximations
unconfirmed at 40\,ms. The full configuration was specified before this
comparison and is not claimed to be the best variant in every task.

For the four MCTS stages, selection scans only admitted actions, while
expansion stores sampled outcomes rather than a full joint transition table.
As stochasticity spreads probability across more successors, repeated visits
to a particular child can become scarce; rollout quality then matters more.
Incremental mean backup uses constant extra storage per edge instead of
storing all previous returns, but rare collision penalties can still make
finite-sample values noisy. These describe computational mechanisms; UCB
selection and the mean-backup rule were not separately ablated.
